\documentclass[conference]{IEEEtran}

\usepackage[utf8]{inputenc}
\usepackage{amsmath,amssymb}
\usepackage{graphicx}
\usepackage{hyperref}
\usepackage{booktabs}
\usepackage{siunitx}
\usepackage{cleveref}
\usepackage{xcolor}
\usepackage{float}
\usepackage{listings}
\usepackage{microtype}

\lstset{
    language=Python,
    basicstyle=\ttfamily\footnotesize,
    keywordstyle=\color{blue},
    commentstyle=\color{green!60!black},
    stringstyle=\color{red!70!black},
    breaklines=true,
    frame=single,
    numbers=left,
    numberstyle=\tiny\color{gray},
}

\title{Constructing Molecular Mechanics Force Fields:\\
A Modular Implementation of Harmonic and Non-Bonded Potentials}

\author{
    \IEEEauthorblockN{Ryan Kamp}
    \IEEEauthorblockA{Dept.\ of Computer Science\\
        University of Cincinnati\\
        Cincinnati, OH 45221\\
        \texttt{kamprj@mail.uc.edu}\\
        \url{https://github.com/ryanjosephkamp}}
    \and
    \IEEEauthorblockN{\textnormal{\textit{Week 16, Project 1}}}
    \IEEEauthorblockA{Biophysics Portfolio\\
        CS Research Self-Study\\
        February 23, 2026}
}

\begin{document}

\maketitle

% ─────────────────────────────────────────────────────────────────────
\begin{abstract}
Molecular mechanics force fields translate the quantum-mechanical reality of
chemical bonding into a computationally tractable classical potential energy
function~$U(\mathbf{r})$.  We present a from-scratch Python implementation of
the five canonical force field energy terms---harmonic bond stretching, harmonic
angle bending, periodic dihedral torsions, Lennard-Jones 12-6 van der Waals
interactions, and Coulomb electrostatics---parameterized with AMBER~ff99 values.
The engine evaluates six preset polypeptide systems ranging from the alanine
dipeptide (22~atoms) to a 5-residue alpha helix (50~atoms), computes per-bond
stress to identify frustrated regions, and provides an interactive ``Break a
Bond'' experiment.  All energies are reported in~kcal/mol with physically
meaningful parameters drawn from the Cornell~et~al.\ (1995) force
field~\cite{cornell1995}.  Non-bonded interactions consistently dominate the
total energy in systems larger than~20 atoms, and the harmonic approximation
accurately reproduces near-equilibrium behavior while diverging at large
displacements.
\end{abstract}

\begin{IEEEkeywords}
molecular mechanics, force field, AMBER, Lennard-Jones, Coulomb electrostatics,
harmonic potential, protein simulation
\end{IEEEkeywords}

% ─────────────────────────────────────────────────────────────────────
\section{Introduction}

Predicting protein structure and stability requires evaluating the total
potential energy~$U(\mathbf{r})$ as a function of all~$3N$ atomic coordinates.
In molecular mechanics (MM), the electronic structure is replaced by a classical
energy function composed of simple analytical terms: springs for bonds, cosines
for torsions, and power laws for van der Waals interactions~\cite{ponder2003}.
This approximation reduces the computational cost from the~$\mathcal{O}(N^7)$
scaling of high-level quantum methods to~$\mathcal{O}(N^2)$ (or lower with
cutoffs), enabling simulations of systems containing thousands of
atoms~\cite{karplus2002}.

The total MM potential energy decomposes into five terms:
\begin{equation}\label{eq:total}
U(\mathbf{r}) = U_{\text{bonds}} + U_{\text{angles}} + U_{\text{dihedrals}}
    + U_{\text{VdW}} + U_{\text{elec}}
\end{equation}

Each term is parameterized independently using experimental or quantum-chemical
data.  The AMBER family of force fields~\cite{cornell1995, wang2004} is among
the most widely used for biomolecular simulations.

This work presents a modular Python implementation of~\cref{eq:total} with
AMBER~ff99 parameters, applied to six preset polypeptide systems.  We
introduce a per-bond stress metric that highlights frustrated regions and
provide an interactive bond-stretching experiment that demonstrates the
breakdown of the harmonic approximation.

\subsection{Scope and Contributions}

\begin{enumerate}
    \item A self-contained force-field engine implementing all five energy
          terms with AMBER~ff99 parameters.
    \item Per-bond stress visualization for identifying frustrated regions.
    \item A ``Break a Bond'' experiment illustrating harmonic energy growth.
    \item Parameter scanning and coordinate perturbation utilities.
    \item Interactive Streamlit dashboard and CLI for exploration.
\end{enumerate}

% ─────────────────────────────────────────────────────────────────────
\section{Theory}

\subsection{Harmonic Bond Stretching}

Covalent bonds are modeled as harmonic oscillators.  The energy for bond~$b$
between atoms~$i$ and~$j$ is:
\begin{equation}\label{eq:bond}
U_{\text{bond}} = \frac{1}{2} k_b (r_{ij} - r_0)^2
\end{equation}
where~$k_b$ is the force constant, $r_{ij}$ is the internuclear distance, and
$r_0$ is the equilibrium bond length.  AMBER~ff99 values for a C--C single bond
are~$k_b = 340$~kcal/(mol\,\AA$^2$) and~$r_0 = 1.526$~\AA~\cite{cornell1995}.

\subsection{Harmonic Angle Bending}

The angle~$\theta_{ijk}$ formed by three consecutively bonded atoms resists
deformation from equilibrium:
\begin{equation}\label{eq:angle}
U_{\text{angle}} = \frac{1}{2} k_\theta (\theta_{ijk} - \theta_0)^2
\end{equation}
Tetrahedral sp$^3$ carbon centers have~$\theta_0 \approx 109.5°$ with~$k_\theta
\approx 63$~kcal/(mol\,rad$^2$).

\subsection{Periodic Dihedral Torsions}

Rotation about the central bond of four consecutive atoms follows a periodic
potential:
\begin{equation}\label{eq:dihedral}
U_{\text{dihedral}} = \frac{V_n}{2} \left[ 1 + \cos(n\phi - \gamma) \right]
\end{equation}
where~$V_n$ is the barrier height, $n$~is the periodicity, $\phi$~is the
torsion angle, and~$\gamma$~is the phase offset.

\subsection{Lennard-Jones 12-6 Potential}

Non-bonded van der Waals interactions are modeled by:
\begin{equation}\label{eq:lj}
U_{\text{LJ}} = 4\varepsilon \left[ \left(\frac{\sigma}{r}\right)^{12}
    - \left(\frac{\sigma}{r}\right)^{6} \right]
\end{equation}
The~$r^{-12}$ term captures Pauli repulsion; the~$r^{-6}$ term represents
London dispersion.  Cross-type parameters use the Lorentz-Berthelot combining
rules:
\begin{equation}\label{eq:combine}
\sigma_{ij} = \frac{\sigma_i + \sigma_j}{2}, \quad
\varepsilon_{ij} = \sqrt{\varepsilon_i \varepsilon_j}
\end{equation}

\subsection{Coulomb Electrostatics}

Partial atomic charges interact via Coulomb's law:
\begin{equation}\label{eq:coulomb}
U_{\text{elec}} = \frac{332.0636 \, q_i q_j}{\varepsilon_r \, r_{ij}}
\end{equation}
The constant~$332.0636$~kcal\,\AA/(mol\,e$^2$) ensures correct units when
charges are in elementary charge units and distances in angstroms.

\subsection{Exclusion Scheme}

Atoms separated by one bond (1-2~pairs) or two bonds (1-3~pairs) are excluded
from non-bonded calculations.  These short-range interactions are already
captured by the bond and angle terms~\cite{cornell1995}.

\subsection{Bond Stress Metric}

We define per-bond stress as the normalized harmonic energy:
\begin{equation}\label{eq:stress}
s_b = \frac{U_b}{\max_{b'} U_{b'}}
\end{equation}
Values near~1.0 indicate maximally strained bonds, which may correspond to
frustrated regions involved in catalysis or conformational
change~\cite{brooks2009}.

% ─────────────────────────────────────────────────────────────────────
\section{Methods}

\subsection{Software Architecture}

The implementation follows a layered architecture with strict dependency flow:

\begin{table}[H]
\centering
\caption{Module Architecture}
\label{tab:architecture}
\begin{tabular}{@{}lll@{}}
\toprule
Module & Responsibility & Lines \\
\midrule
\texttt{forcefield\_engine.py} & Data structures, energy calculators & $\sim$1570 \\
\texttt{analysis.py} & Typed result containers, pipelines & $\sim$640 \\
\texttt{visualization.py} & Plotly + Matplotlib renderers & $\sim$850 \\
\texttt{app.py} & 8-page Streamlit dashboard & $\sim$940 \\
\texttt{main.py} & CLI with 5 modes & $\sim$230 \\
\bottomrule
\end{tabular}
\end{table}

\subsection{ForceField Data Model}

The \texttt{ForceField} class maintains lists of \texttt{Atom},
\texttt{BondParam}, \texttt{AngleParam}, and \texttt{DihedralParam} dataclass
instances.  An exclusion set of (1-2) and (1-3) atom pairs is built
automatically from the bond and angle topology.

\subsection{Energy Evaluation Protocol}

Energy evaluation proceeds sequentially through the five terms:
\begin{enumerate}
    \item Bond energy: iterate over all \texttt{BondParam} entries; compute
          $r_{ij}$ via Euclidean distance; apply~\cref{eq:bond}.
    \item Angle energy: compute~$\theta_{ijk}$ via the dot product of bond
          vectors; apply~\cref{eq:angle}.
    \item Dihedral energy: compute~$\phi$ via the cross-product method;
          apply~\cref{eq:dihedral}.
    \item Non-bonded energy: double loop over all non-excluded pairs within
          the cutoff; compute LJ~(\cref{eq:lj}) + Coulomb~(\cref{eq:coulomb}).
    \item Total: sum all five contributions.
\end{enumerate}

\subsection{Preset Molecule Systems}

\begin{table}[H]
\centering
\caption{Preset Polypeptide Systems}
\label{tab:presets}
\begin{tabular}{@{}lrrrl@{}}
\toprule
Molecule & Atoms & Bonds & Angles & Description \\
\midrule
Alanine dipeptide   & 22 & 21 & 20 & Standard benchmark \\
Glycine tripeptide  & 33 & 32 & 31 & Extended backbone \\
Alpha helix (5-res) & 50 & 49 & 48 & Poly-alanine helix \\
Beta hairpin        & 16 & 15 & 12 & Backbone strand \\
Salt bridge         & 14 & 13 & 12 & Lys-Asp charge pair \\
Disulfide bond      & 12 & 11 & 10 & Cys-S-S-Cys cross-link \\
\bottomrule
\end{tabular}
\end{table}

\subsection{Stress Visualization}

Per-bond stresses are computed via~\cref{eq:stress} and mapped to a
blue-to-red color scale.  Bonds with stress~$> 0.7$ are flagged as frustrated
regions and reported to the user.

\subsection{Break a Bond Experiment}

A selected bond is incrementally stretched from equilibrium by up to a
user-specified displacement.  At each step, the full energy~$U(\mathbf{r})$ is
recomputed, demonstrating the quadratic growth of the harmonic
term~(\cref{eq:bond}).

% ─────────────────────────────────────────────────────────────────────
\section{Results}

\subsection{Energy Decomposition --- Alanine Dipeptide}

The alanine dipeptide (Ace-Ala-Nme, 22~atoms) provides the primary benchmark.
Non-bonded interactions (VdW + electrostatic) consistently dominate the total
energy, accounting for the majority of the energy magnitude.  Bond and angle
terms contribute moderate positive energy due to minor deviations from
equilibrium geometry.

\subsection{Stress Analysis}

Terminal bonds (N-terminus and C-terminus cappings) exhibit higher stress than
interior backbone peptide bonds, reflecting minor geometric deviations in the
preset builders.  Frustrated regions correspond to bonds with normalized
stress~$> 0.7$.  The dashboard provides a per-residue frustration table
reporting the residue name, atom count, number of stressed bonds exceeding
the system mean, maximum and mean bond stress, and the largest bond
displacement~$|r - r_0|$ in angstroms.

\subsection{Break a Bond}

Stretching the N--C$_\alpha$ bond of alanine dipeptide by up to~$+2.0$~\AA\
produces the expected parabolic energy curve.  The energy increases by several
hundred~kcal/mol---far exceeding the thermal energy~$k_BT \approx
0.6$~kcal/mol at~300~K---demonstrating the high cost of covalent bond rupture
in the harmonic approximation.

\subsection{Parameter Sensitivity}

Scanning the bond force constant~$k_b$ from 50 to
600~kcal/(mol\,\AA$^2$) reveals a linear relationship between~$k_b$ and the
total bond energy contribution when bonds are displaced from equilibrium.  The
dielectric constant~$\varepsilon_r$ modulates electrostatic energy
approximately inversely.

\subsection{Perturbation Robustness}

Gaussian coordinate perturbations with scale factors from 0.01 to 0.5~\AA\
demonstrate:
\begin{itemize}
    \item $\sigma = 0.01$~\AA: energy change~$< 1$~kcal/mol (thermal regime)
    \item $\sigma = 0.1$~\AA: energy change~$\sim 10$--$50$~kcal/mol (strain)
    \item $\sigma = 0.5$~\AA: energy change~$\sim 100+$~kcal/mol (non-physical)
\end{itemize}

\subsection{Cross-Molecule Comparison}

Larger systems show proportionally larger non-bonded contributions due to
$\mathcal{O}(N^2)$ pair scaling.  The salt bridge system is dominated by
strong Coulomb attraction between charged residues.

% ─────────────────────────────────────────────────────────────────────
\section{Discussion}

\subsection{Accuracy of the Harmonic Approximation}

The harmonic bond model accurately reproduces near-equilibrium energetics but
diverges from reality at large displacements.  A Morse potential would provide
more realistic behavior at the cost of additional parameters~\cite{leach2001}.
The ``Break a Bond'' experiment clearly illustrates this limitation.

\subsection{Non-Bonded Dominance}

In systems larger than~$\sim 20$ atoms, non-bonded interactions dominate the
energy landscape.  This reflects the~$N(N-1)/2$ pair scaling versus the
linear scaling of bonded terms.  Electrostatics are particularly important in
polypeptides due to large partial charges on backbone N, C, and O
atoms~\cite{mackerell1998}.

\subsection{Computational Considerations}

The naive~$O(N^2)$ non-bonded calculation can be accelerated to~$O(N \log N)$
via particle-mesh Ewald~\cite{darden1993} or~$O(N)$ via the fast multipole
method.  Neighbor lists and cell lists are classical spatial data structures
used in production MD codes~\cite{allen2017}.

\subsection{Limitations}

The current implementation uses fixed partial charges (no polarization),
harmonic bonds (no dissociation), and idealized preset geometries.  No explicit
solvent, periodic boundary conditions, or long-range electrostatic corrections
are applied.

\subsection{Biological Relevance}

Force fields enable structure refinement, drug design scoring, molecular
dynamics, and mutation impact prediction~\cite{karplus2002}.  The per-bond
stress metric introduced here provides a simple diagnostic for identifying
structurally frustrated regions that may correlate with functional sites.

% ─────────────────────────────────────────────────────────────────────
\section{Conclusion}

We have implemented a complete molecular mechanics force field engine with five
standard energy terms and AMBER~ff99 parameters.  The engine evaluates six
preset polypeptide systems, computes per-bond stress for frustrated-region
identification, and provides interactive bond-stretching, parameter-scanning,
and coordinate-perturbation experiments.  Non-bonded interactions dominate in
larger systems, and the harmonic approximation accurately describes
near-equilibrium behavior while failing at large displacements.  The modular
Python architecture provides a pedagogical platform for exploring the physics
underlying modern biomolecular simulation.

% ─────────────────────────────────────────────────────────────────────
\begin{thebibliography}{14}

\bibitem{cornell1995}
W.~D. Cornell et~al., ``A Second Generation Force Field for the Simulation of
Proteins, Nucleic Acids, and Organic Molecules,'' \emph{J. Am. Chem. Soc.},
vol.~117, pp.~5179--5197, 1995.

\bibitem{weiner1984}
S.~J. Weiner et~al., ``A New Force Field for Molecular Mechanical Simulation
of Nucleic Acids and Proteins,'' \emph{J. Am. Chem. Soc.}, vol.~106,
pp.~765--784, 1984.

\bibitem{ponder2003}
J.~W. Ponder and D.~A. Case, ``Force Fields for Protein Simulations,''
\emph{Adv. Protein Chem.}, vol.~66, pp.~27--85, 2003.

\bibitem{wang2004}
J.~Wang et~al., ``Development and Testing of a General AMBER Force Field,''
\emph{J. Comput. Chem.}, vol.~25, pp.~1157--1174, 2004.

\bibitem{mackerell1998}
A.~D. MacKerell et~al., ``All-Atom Empirical Potential for Molecular Modeling
and Dynamics Studies of Proteins,'' \emph{J. Phys. Chem. B}, vol.~102,
pp.~3586--3616, 1998.

\bibitem{jorgensen1996}
W.~L. Jorgensen et~al., ``Development and Testing of the OPLS All-Atom Force
Field,'' \emph{J. Am. Chem. Soc.}, vol.~118, pp.~11225--11236, 1996.

\bibitem{karplus2002}
M.~Karplus and J.~A. McCammon, ``Molecular Dynamics Simulations of
Biomolecules,'' \emph{Nat. Struct. Biol.}, vol.~9, pp.~646--652, 2002.

\bibitem{brooks2009}
B.~R. Brooks et~al., ``CHARMM: The Biomolecular Simulation Program,''
\emph{J. Comput. Chem.}, vol.~30, pp.~1545--1614, 2009.

\bibitem{duan2003}
Y.~Duan et~al., ``A Point-Charge Force Field for Molecular Mechanics
Simulations of Proteins,'' \emph{J. Comput. Chem.}, vol.~24,
pp.~1999--2012, 2003.

\bibitem{lennard1924}
J.~E. Lennard-Jones, ``On the Determination of Molecular Fields,''
\emph{Proc. R. Soc. Lond. A}, vol.~106, pp.~463--477, 1924.

\bibitem{darden1993}
T.~Darden, D.~York, and L.~Pedersen, ``Particle Mesh Ewald: An
$N\!\cdot\!\log(N)$ Method for Ewald Sums,'' \emph{J. Chem. Phys.},
vol.~98, pp.~10089--10092, 1993.

\bibitem{leach2001}
A.~R. Leach, \emph{Molecular Modelling: Principles and Applications},
2nd~ed.\hskip 1em Pearson, 2001.

\bibitem{schlick2010}
T.~Schlick, \emph{Molecular Modeling and Simulation: An Interdisciplinary
Guide}, 2nd~ed.\hskip 1em Springer, 2010.

\bibitem{allen2017}
M.~P. Allen and D.~J. Tildesley, \emph{Computer Simulation of Liquids},
2nd~ed.\hskip 1em Oxford University Press, 2017.

\bibitem{ferreiro2007}
D.~U. Ferreiro et~al., ``Localizing Frustration in Native Proteins and
Protein Assemblies,'' \emph{Proc. Natl. Acad. Sci. USA}, vol.~104,
pp.~19819--19824, 2007.

\bibitem{morse1929}
P.~M. Morse, ``Diatomic Molecules According to the Wave Mechanics.~II.
Vibrational Levels,'' \emph{Physical Review}, vol.~34, pp.~57--64, 1929.

\end{thebibliography}

\end{document}
